\documentclass{article}
\usepackage{graphicx} % Required for inserting images
\usepackage{amsmath}
\usepackage[a4paper,margin=1.5cm]{geometry}

\title{ProjetIntegrateurMPC}
\author{mathis.poteau28 }
\date{December 2025}

\begin{document}

\maketitle

\section{Introduction}



\section{Model}

On considère un réseau orienté $G=(V,E)$ composé de deux types de sommets :
\begin{itemize}
    \item $\mathcal{T}$ : ensemble des \textbf{cuves} (réservoirs),
    \item $\mathcal{X}$ : ensemble des \textbf{cellules de transport} (tunnel).
\end{itemize}
On a $V = \mathcal{T} \cup \mathcal{X}$ et $\mathcal{T} \cap \mathcal{X} = \emptyset$.

Le temps est discrétisé avec un pas $h > 0$ :
\[
t_k = k h, \quad k = 0,1,\dots,K.
\]

\subsection{Variables d'état et de commande}

Pour tout $k$ :
\begin{itemize}
    \item $x_T(k)$ : volume d’eau dans la cuve $T \in \mathcal{T}$,
    \item $x_i(k)$ : volume d’eau dans la cellule $i \in \mathcal{X}$,
    \item $u_T(k) \in [0,1]$ : commande d’ouverture de la vanne de la cuve $T$,
    \item $f_{ij}(k) \ge 0$ : flux entre les sommets $i \to j$,
    \item $\mu_i(k)$ : débit de sortie hors réseau pour les cellules terminales.
\end{itemize}

\subsection{Fonctions de demande et d’offre}

Pour chaque cellule $i \in \mathcal{X}$, on définit :

\paragraph{Fonction de demande (d’après le code Python)}
\[
d_i(x_i(k)) =\gamma_i  \min\left(  \frac{v_i}{l_i} x_i(k) ,\; C_i \right),
\]
où $v_i$ est la vitesse, $l_i$ la longueur de la cellule, $C_i$ la capacité maximale et $\gamma_i \in \begin{bmatrix}
    0, 1
\end{bmatrix}$ une commande permettant de contrôler la vitesse des véhicules.

\paragraph{Fonction d’offre (supply)}
\[
s_i(x_i(k)) = \max(c_i - w_i x_i(k) ,\; 0),
\]
où $c_i$ est une constante de capacité et $w_i$ la pente de saturation.

\subsection{Définition des flux}

\paragraph{Flux cellule $\to$ cellule}
\[
f_{ij}(k) = \min\left(d_i(x_i(k)),\; s_j(x_j(k))\right), \quad i,j \in \mathcal{X}.
\]

\paragraph{Flux cuve $\to$ cellule}
\[
f_{Tj}(k) = \min\left(u_T(k) \, x_T(k),\; s_j(x_j(k))\right), 
\quad T \in \mathcal{T},\ j \in \mathcal{X}.
\]

\paragraph{Sortie hors réseau}
Si une cellule $i$ n’a pas de successeur :
\[
\mu_i(k) = d_i(x_i(k)).
\]

\paragraph{Interdiction de retour vers les cuves}
\[
f_{jT}(k) = 0, \quad \forall j \in \mathcal{X},\ \forall T \in \mathcal{T}.
\]

\subsection{Dynamiques discrètes du système}

\paragraph{Évolution des cuves}
\[
x_T(k+1) = x_T(k) 
- h \sum_{j:\,T \to j} f_{Tj}(k).
\]

\paragraph{Évolution des cellules}
\[
x_i(k+1)
= x_i(k)
+ h \left(
\sum_{j:\,j \to i} f_{ji}(k)
- \sum_{m:\,i \to m} f_{im}(k)
- \mu_i(k)
\right),
\quad i \in \mathcal{X}.
\]

\subsection{Contraintes physiques}

\[
\begin{cases}
x_i(k) \ge 0, \quad x_T(k) \ge 0, \\
f_{ij}(k) \ge 0, \\
0 \le u_T(k) \le x_T(k).
\end{cases}
\]

\subsection{Relaxation convexe}

Afin de rendre le problème convexe, les flux $f_{ij}(k)$ sont considérés comme variables indépendantes et les fonctions $\min(\cdot)$ sont remplacées par des contraintes d'inégalité :

\[
0 \le \bar{f}_{ij}(k) \le d_i(x_i(k)), \quad
\sum_{j:\,j\to i} \bar{f}_{ji}(k) \le s_i(x_i(k)).
\]

\subsection{Problème de contrôle optimal sur horizon fini}

On considère le critère quadratique :
\[
J = \sum_{k=0}^{K} \left( \sum_{i \in \mathcal{X}} x_i(k) + \sum_{T \in \mathcal{T}} x_T(k) \right)
\]


sous contraintes :
\begin{itemize}
    \item contraintes physiques sur la commande : 
    $\begin{cases}
        0 \le u_T(k) \le x_T(k)$ for $T \in \mathcal{T}, k \in 
    \begin{bmatrix}
        0, K
    \end{bmatrix} \\
    \gamma_i \in 
    \begin{bmatrix}
        0, 1
    \end{bmatrix}
    \end{cases}$
    
    \item contraintes de flux :
\[
\sum_{m:\,i \to m}  \bar{f}_{im}(k) \le d_i(x_i(k)) ,
\]
\[
\sum_{j:\,j \to i} \bar{f}_{ji}(k) \le s_i(x_i(k)),
\]
\[
\sum_{j:\,T \to j} f_{Tj}(k) \le u_T(k) x_T(k)
\]
    \item dynamiques du système : 
\[
x_i(k+1)
= x_i(k)
+ h \left(
\sum_{j:\,j \to i} \bar{f}_{ji}(k)
- \sum_{m:\,i \to m} \bar{f}_{im}(k)
- \mu_i(k)
\right),
\quad i \in \mathcal{X}.
\]
    \item pas de retour vers les cuves :
\[
f_{jT}(k) = 0,
\]
    \item conditions initiales :
\[
x(0) \text{ donné.}
\]
\end{itemize}





\end{document}
